% Source and dependency audit for the exact P09.2 reduction packet.

\paragraph{Scope and status.}
\Cref{thm:p09p2-reduction} is a source-derived partial reduction.  No
novelty or priority claim is made.  It does not prove that an ample symmetric
cubical metric in non-isotrivial mixed reduction is an ordinary
delta-metric, and it gives no counterexample.  Thus it does not close P09.2.
Its exact output is: the \(\operatorname{Pic}^0\) part is unconditional in
the stated GK field setting; the remaining question on a fixed \(A\) is
equivalent to the ample symmetric case.

\paragraph{Primary-source locations and consumers.}
\begin{enumerate}
\item Gubler--K\"unnemann,
\emph{A tropical approach to non-Archimedean Arakelov geometry},
\cite{GK17}:
 \begin{enumerate}
 \item Proposition~8.11 supplies formal metric
 \(\Longleftrightarrow\) piecewise linear metric;
 \item Proposition~8.13 supplies algebraic metric
 \(\Longleftrightarrow\) formal metric on a proper variety;
 \item Example~8.16 supplies an \(N\geq1\) for which the normalized
 canonical metric of an algebraically trivial line bundle has algebraic
 \(N\)-th tensor power;
 \item Definition~9.9 and \S9.13 fix the ordinary tropical-chart meaning
 of delta-metric and the roman differential convention
 \(\mathrm d'_{\mathrm P}\mathrm d''_{\mathrm P}\);
 \item \S9.14 supplies locality, isometry invariance, tensor and dual
 closure, and the positive-tensor-power root implication
 \eqref{eq:p09p2-root-rule};
 \item Remark~9.16 supplies piecewise linear/formal/algebraic
 \(\Longrightarrow\) delta-metric.
 \end{enumerate}
These statements consume every arrow in the recognition chain and every
closure step in \cref{itm:p09p2-pic0,itm:p09p2-split}.

\item Gubler, \emph{Non-Archimedean canonical measures on abelian
varieties}, \cite[Example~3.7]{Gub10CM}, is the algebraicity input cited by
\cite[Example~8.16]{GK17}.  The proof packet invokes it only through that
exact GK statement.

\item Chambert--Loir--Ducros, version of 25 July 2025,
\cite[Lemmas~19.1.4 and 19.1.6]{CLD25}, supplies existence and uniqueness
after rigidification of cubical metrics, compatibility with tensor product,
duality, and homomorphism pullback, and the exponent-one characterization
for antisymmetric line bundles.  The direct calculation from
\eqref{eq:p09p2-d3} separately verifies the particular compatibility used
in \cref{itm:p09p2-split}.

\item Milne, \emph{Abelian Varieties}, version 2.00,
\cite[Proposition~8.4, Remark~8.5, and Remark~10.19]{milneAV}, supplies the
characterization of \(\operatorname{Pic}^0(A)\), inversion-as-duality on
\(\operatorname{Pic}^0(A)\), and the identity action of inversion on
\(\operatorname{NS}(A)\).  These statements are used, respectively, to
match the normalized cubical metric to the exponent-one canonical metric
and to prove \(L_-\in\operatorname{Pic}^0(A)\).  The cited version is at
\url{https://www.jmilne.org/math/CourseNotes/AV.pdf}.

\item The Stacks Project, \cite{stacks-project}: Lemma~39.9.2 supplies an
ample line bundle on \(A\); Lemma~28.27.5 (Tag~0890) supplies ampleness of
the symmetric tensor product; and Lemma~29.38.7 (Tag~0892) supplies
ampleness of \eqref{eq:p09p2-ample-twist} for \(m\gg0\).  These are exactly
the algebraic inputs consumed by \cref{itm:p09p2-ample-reduction}.
\end{enumerate}

\paragraph{No-jump audit.}
There is no base-change or descent step: the field hypotheses are exactly
those imposed throughout \cite{GK17}.  The proof never infers ordinary
delta-metricity from weak smoothness.  The positive-power step is the
explicit root closure in \cite[\S9.14]{GK17}, not an extraction of a model
root.  The isometry \eqref{eq:p09p2-square-split} uses induced metrics and
therefore carries no hidden scalar.  Arbitrary normalization is handled by
a constant smooth metric.  Finally, \eqref{eq:p09p2-ample-twist} changes
only the symmetric line bundle \(S\); the metric used on it is the induced
cubical tensor metric, so the hypothesis in
\cref{itm:p09p2-ample-symmetric} applies without a uniqueness rescaling.
